\chapter{Split algebra and exact factorization of the character}
\label{chap:factorizacion-split}
\index{split algebra}\index{character!factorization split}
\index{null cone}\index{channels $q_\pm$}

The formal signature of an extension and the angular channels follow from branches already constructed. In this section, both converge in a determinant factorization that separately preserves the radial scale and the relative orientation. The result belongs to the algebraic stratum: it does not by itself identify the null ideals with physical photons or the positive interior with a material particle.

\section{Conjugate angular coordinates and bidirectional spectral lattice}
\label{sec:reticula-angular-split}

The preceding causal chain fixes
\[
 A_{\rm rad}=\frac{\pi A}{180},\qquad
 \Cstarrad=\frac{\pi\Cstar}{180},
\]
and, only after obtaining the direct coordinate \(A\) and the unique
conjugate coordinate \(C^*\),
\[
 q_+=\exp[-(A_{\rm rad}+\Cstarrad)],
 \qquad
 q_-=\exp[-(A_{\rm rad}-\Cstarrad)].
\]
For an angular signature \((n_A,s_C)\), we introduce the two winding
numbers
\[
 W_+=6n_A+s_C,
 \qquad
 W_-=6n_A-s_C.
\]

\begin{theorem}[Angular lattice of index twelve]
\label{thm:reticula-two-way}
The map
\[
 \mathcal W:\ZZ^2\to\ZZ^2,
 \qquad
 \binom{n_A}{s_C}\longmapsto
 \binom{W_+}{W_-}
 =\begin{pmatrix}6&1\\6&-1\end{pmatrix}
 \binom{n_A}{s_C}
\]
is injective, and its image has index twelve. Its invariant factors are
\((1,12)\), and
\[
 \operatorname{im}\mathcal W
 =\{(u,v)\in\ZZ^2:u+v\equiv0\pmod{12}\}.
\]
Moreover,
\[
 \boxed{
 \exp\!\left(n_AA_{\rm rad}+\frac{s_C}{6}\Cstarrad\right)
 =\left(q_+^{-W_+}q_-^{-W_-}\right)^{1/12}.}
\]
\end{theorem}

\begin{proof}
The matrix has determinant \(-12\). Since the greatest common divisor of its
entries is one, its Smith normal form is
\(\operatorname{diag}(1,12)\). If
\((u,v)=(6n+s,6n-s)\), then \(u+v=12n\). Conversely, if
\(u+v\) is a multiple of twelve, then
\[
 n=\frac{u+v}{12},
 \qquad
 s=\frac{u-v}{2}
\]
are integers and recover \((u,v)\). Finally,
\[
 \frac{W_+(A_{\rm rad}+\Cstarrad)
       +W_-(A_{\rm rad}-\Cstarrad)}{12}
 =n_AA_{\rm rad}+\frac{s_C}{6}\Cstarrad,
\]
which proves the identity upon exponentiation.
\end{proof}

The involution \(s_C\mapsto-s_C\) interchanges
\(W_+\leftrightarrow W_-\) and \(q_+\leftrightarrow q_-\). The twelfth
root is the integral index of the image of \(\mathcal W\), not a
rounding. The coordinate \(C^*\) acts as the antisymmetric component of the lattice;
it is neither a decimal correction to a mass nor an output of
\(\Gamma_{6\to8}\).

\section{Real split algebra}
\label{sec:algebra-real-partida}

\begin{definition}[Real split algebra]
Let
\[
 \mathcal A_{\rm sp}
 :=\RR[\mathsf R]/(\mathsf R^2-1).
\]
Conjugation and the norm are defined, respectively, by
\[
 \overline{a+b\mathsf R}=a-b\mathsf R,
 \qquad
 N_{\rm sp}(z)=z\overline z=a^2-b^2.
\]
\end{definition}

Conjugate idempotents
\[
 P_\pm=\frac{1\pm\mathsf R}{2}
\]
satisfy
\[
 P_\pm^2=P_\pm,\qquad
 P_+P_-=0,\qquad
 P_++P_-=1,\qquad
 \overline{P_+}=P_-.
\]

\begin{proposition}[Split spectral decomposition]
\label{prop:descomposicion-espectral-partida}
Every \(z\in\mathcal A_{\rm sp}\) admits a unique decomposition
\[
 z=r_+P_++r_-P_-,
 \qquad r_\pm\in\RR,
\]
and its norm is
\[
 \boxed{N_{\rm sp}(z)=r_+r_-.}
\]
Each principal ideal \(\mathcal A_{\rm sp}P_\pm=\RR P_\pm\) is a null
line; in particular,
\[
 N_{\rm sp}(\lambda P_\pm)=0.
\]
\end{proposition}

\begin{proof}
If \(z=a+b\mathsf R\), the choice \(r_+=a+b\), \(r_-=a-b\) provides the
decomposition, whose uniqueness follows from the independence of
\(P_+\) and \(P_-\). Conjugation exchanges the two coefficients:
\(\overline z=r_-P_++r_+P_-\). The orthogonality of the idempotents gives
\[
 z\overline z=(r_+r_-)(P_++P_-)=r_+r_-.
\]
The statement about each ideal is obtained by nullifying one of the two coefficients.
\end{proof}

\begin{definition}[Causal cone partitioned]
\label{def:cono-causal-partido}
The positive cone, its null boundary, and its interior are
\[
 \mathcal C_{\rm sp}^{+}
 =\{r_+P_++r_-P_-:r_+,r_-\geq0\},
\]
\[
 \partial_{\rm null}\mathcal C_{\rm sp}^{+}
 =\RR_{\geq0}P_+\cup\RR_{\geq0}P_-,
\qquad
 \operatorname{int}\mathcal C_{\rm sp}^{+}
 =\{r_+P_++r_-P_-:r_+r_->0\}.
\]
\end{definition}

The boundary has a single active direction and zero norm. The interior requires
two simultaneous directions and positive norm. This stratification is exact;
its physical realization requires maps that preserve the pertinent norm, transport, and
dynamics.

\section{Relative transporter and representational realization conditions}
\label{sec:transportador-relativo-split-md}

The historical block \(B_0=7I+2\mathsf R\) has spectrum \(\{9,5\}\), and the
angular block \(B_1=AI+\Cstar\mathsf R\) has spectrum
\(\{A+\Cstar,A-\Cstar\}\). The
\cref{thm:no-go-entrelazador-bloques-split} proves that, with the current internal
values, every putative linear intertwiner
\(\Phi B_0=B_1\Phi\) is zero. The active object is the relative
multiplicative transporter already constructed, not an equivalence between
representations with disjoint spectra.

\section{Partial factorization of the angular character}
\label{sec:factorizacion-split-angular}

Let
\[
 \sigma=\frac{s_C}{6},\qquad
 w_\pm=\frac{n_A\pm\sigma}{2}=\frac{W_\pm}{12},
 \qquad
 r_\pm=q_\pm^{-w_\pm},
\]
and let us define
\[
 Z_v=r_+P_++r_-P_-.
\]

\begin{theorem}[Partial factorization of the angular character]
\label{thm:split-character}
For every angular signature \(v=(n_A,s_C)\in\ZZ^2\),
\[
 \boxed{
 N_{\rm sp}(Z_v)
 =\det Z_v
 =\exp\!\left(
 n_AA_{\rm rad}+\frac{s_C}{6}\Cstarrad
 \right).}
\]
\end{theorem}

\begin{proof}
By the definition of the channels,
\[
 \log r_+
 =\frac{n_A+\sigma}{2}(A_{\rm rad}+\Cstarrad),
 \qquad
 \log r_-
 =\frac{n_A-\sigma}{2}(A_{\rm rad}-\Cstarrad).
\]
Upon summation,
\[
 \log(r_+r_-)
 =n_AA_{\rm rad}+\sigma\Cstarrad.
\]
The \cref{prop:descomposicion-espectral-partida} gives
\(N_{\rm sp}(Z_v)=r_+r_-\), and the diagonal representation in the basis
\((P_+,P_-)\) gives the same expression for the determinant.
\end{proof}

The component not seen by the determinant preserves the relative orientation:
\[
 \frac12\log\frac{r_+}{r_-}
 =\frac12\left(
 n_A\Cstarrad+\frac{s_C}{6}A_{\rm rad}
 \right),
\]
\[
 Z_v=
 \exp\!\left[
 \frac12\left(n_AA_{\rm rad}+\frac{s_C}{6}\Cstarrad\right)
 \right]
 \exp\!\left[
 \frac12\left(n_A\Cstarrad+\frac{s_C}{6}A_{\rm rad}\right)
 \mathsf R
 \right].
\]
The first factor fixes the radial scale; the second records the relative
orientation. They are two coordinated readings of one and the same element.

\section{Central incorporation of the towers}
\label{sec:factorizacion-split-torres}

Once \(\Dfour\), \(s_{120}\), and
\(\zCorr=\zeta_{270}^{\rm corr}\) have been constructed globally, let
\[
 t(v)=
 k_{\Dfour}\Dfour
 +\nu_{120}s_{120}
 +\nu_{270}\zCorr,
 \qquad
 \widetilde Z_v=e^{t(v)/2}Z_v.
\]

\begin{corollary}[Complete character as original norm]
\label{cor:split-character-torres}
Having fixed the preceding global blocks,
\[
 \boxed{
 N_{\rm sp}(\widetilde Z_v)
 =\exp\!\left(
 n_AA_{\rm rad}+\frac{s_C}{6}\Cstarrad
 +k_{\Dfour}\Dfour+\nu_{120}s_{120}
 +\nu_{270}\zCorr
 \right).}
\]
\end{corollary}

\begin{proof}
In the two-dimensional representation, multiplication by the central scalar
\(e^{t(v)/2}\) multiplies the determinant by \(e^{t(v)}\). The
\cref{thm:split-character} is applied.
\end{proof}

The corollary constructs the declared split realization of the active character;
it does not classify all possible realizations or transform the towers into
independent proofs. Its hypotheses are the already fixed channels and global
blocks, and its negative control is immediate: omitting the central factor
\(1/2\) doubles the tower exponent; simultaneously activating both
idempotents with positive coefficients departs from the null boundary.
